\documentclass[10pt,a4paper]{amsart}
\usepackage[margin=19mm]{geometry}
\usepackage{amsmath,amssymb,mathtools}
\usepackage[scr=rsfs,cal=euler]{mathalfa}
\usepackage{xfrac}
\usepackage[unicode,hidelinks,bookmarks=false]{hyperref}
\newcommand{\ie}{i.e.\ }
\newcommand{\cf}{cf.\ }
\newcommand{\resp}{respectively\ }
\newcommand{\Subsection}{\subsection}
\providecommand{\nobreakdash}{\nobreak\hbox{-}}
\setlength{\emergencystretch}{2em}
\multlinegap0pt
\usepackage{bm}
\usepackage{bbm}
\usepackage{dsfont}
\usepackage{xspace}
\usepackage{cancel}
\usepackage{mathtools}
\usepackage{amsthm}
\makeatletter
\@ifundefined{prop}{\newtheorem{prop}{Proposition}}{}
\makeatother
\newcommand{\C}{\mathbb{C}}
\newcommand{\R}{\mathbb{R}}
\title[On partial diffusion and mixing without hypoellipticity]{On partial diffusion and mixing without hypoellipticity}
\author{Xu'an Dou, Delphine Salort, Didier Smets}
\date{Source: arXiv:2511.05280v1 (CC BY 4.0)}
\begin{document}
\maketitle

\noindent\emph{Source and licence.} On partial diffusion and mixing without hypoellipticity. Version arXiv:2511.05280v1 (CC BY 4.0), distributed under the Creative Commons Attribution 4.0 International licence, as recorded in the arXiv licence metadata for this exact version and shown on the abstract page https://arxiv.org/abs/2511.05280v1. Full rights notice: licenses/arxiv-2511.05280-CC-BY-4.0.txt.\par\smallskip

\noindent\textit{Editorial scope.} Counted unit: the Prop whose source label is \texttt{prop:v1} in the article (source0.tex lines 873--880, source label \texttt{prop:v1}), delivered together with the article's own complete original proof of it (source0.tex lines 882--1083, one \texttt{proof} environment of 202 lines). No mathematics is written, repaired or re-derived: statement and proof are the article's own text line for line. Required context retained uncounted: eq:defomega2 (lines 868--871). Every definition, display and label of those lines is retained unchanged. Omitted from this package and not redistributed: 1--867, 872--872, 881--881, 1084--1881. Nothing inside a retained excerpt is omitted or altered: every retained body is delivered exactly as the article's lines are written, so no omission receipt of any kind accompanies this package. Citations: the retained excerpts carry no citation command at all, so no bibliography entry of the article is transcribed and no reference is redistributed. Reference adaptation: none was needed and none was made; every reference inside the delivered unit resolves inside it, so no unresolved reference or citation is produced. Printed slips kept literal: no printed word, symbol, hypothesis or number of the counted statement or of its proof was altered. Layout and class: the article's own class is \texttt{article} with packages including fontenc, babel, geometry, amsmath, amssymb, amsthm, bbm, hyperref, graphicx, xcolor; the wrapper is the lane's pinned \texttt{amsart} editorial wrapper at 10pt on A4, which restates the theorem-like environments the retained text needs and adds the article's own macro definitions that the retained text needs (\texttt{\textbackslash C}, \texttt{\textbackslash R}), with extra wrapper packages (bm, bbm, dsfont, xspace, cancel, mathtools). The delivered numbering is the unit's own. Assets: no retained span contains a figure environment, a graphics-inclusion command, a PDF-page inclusion, a table or a TikZ picture; no image file is copied into this package, the package declares no asset dependency and no image file is redistributed with it. Licence: version arXiv:2511.05280v1 of this article is distributed under the Creative Commons Attribution 4.0 International licence, as recorded in the arXiv licence metadata for this exact version and shown on the abstract page https://arxiv.org/abs/2511.05280v1. A full rights notice is delivered at licenses/arxiv-2511.05280-CC-BY-4.0.txt. Characters: every retained excerpt is pure ASCII, so the delivered engine, XeLaTeX with its default Latin Modern text font, reports no missing character.
\begin{equation}\label{eq:defomega2}
	\omega_2(V) := \max \left\{\int_I V\varphi e_1^2,\ \varphi \in L,
	\|\varphi\|_L \leq 1\right\}.
\end{equation}

\begin{prop}\label{prop:v1} 
Let $V \in L^\infty(I,\R)$ be non constant. Then $\omega_2(V) > 0$ and we have  
\begin{equation}\label{eq:prop1}
r(\lambda_1, V) \geq \frac{\omega_2(V)^2}{18} 
\left( \frac{\pi^2}{|I|^2} + \frac{|I|^2}{\pi^2} {\it Osc}(V)^2\right)^{-1} > 0,
\end{equation}
where ${\it Osc}(V) := {\it ess\,sup}(V) - {\it ess\,inf}(V)$.
\end{prop}

\begin{proof}
The fact that $\omega_2(V) > 0$ when $V$ is non constant is immediate.
Indeed, by the du Bois-Reymond lemma there exists 
$\phi \in \mathcal{D}(I)$ such that $\int_I \phi = 0$ and $\int_I
V \phi \neq 0.$ Since $e_1$ is smooth and positive in the interior of
$I$, we may set $\varphi := \phi / e_1^2$, and we have $\varphi \in
\mathcal{D}(I)$, $\int_I \varphi e_1^2 = 0$, and $\int_I
V \varphi e_1^2 \neq 0.$ The conclusion follows by multiplying
$\varphi$ by a suitable non zero constant. 

We now turn to the estimate of $r(\lambda_1).$ We write any 
$z \in \C$ s.t. ${\rm Re}(z) = \lambda_1$  as 
$z = \lambda_1 + is$, with $s \in \R.$ For $f \in L^2(I)$, 
we wish to estimate a solution $u \in D(A)$ of the equation 
\begin{equation}\label{eq:resolv}
-\partial_{xx}u - \lambda_1 u + i(V - s)u = f.
\end{equation}
For that purpose, we decompose $u$ as 
\begin{equation}\label{eq:decomp-u}
    u = \langle u, e_1\rangle e_1 + v,
\end{equation} and let $C_1 := \langle u, e_1 \rangle.$ Since $v$ is orthogonal to $e_1$ we have
\begin{equation}\label{eq:lambda2}
\int_I |\partial_x v|^2 \geq \lambda_2 \int_I |v|^2,
\end{equation}
where recall $\lambda_2>\lambda_1$ is the second eigenvalue of the corresponding 
\textit{Laplace} operator. Therefore
\begin{equation}\label{eq:spectralgap1}
\int_I |\partial_x v|^2 - \lambda_1 |v|^2 \geq \big(1 -
\frac{\lambda_1}{\lambda_2}\big)
\int_I |\partial_x v|^2 \geq (\lambda_2 - \lambda_1) \int_I
|v|^2.
\end{equation}
Multiplying equation \eqref{eq:resolv} by the conjugate $\bar u$ and then retaining the real
part, we obtain 
$$
\int_I |\partial_x u|^2 - \lambda_1 |u|^2 \leq \|f\|\: \|u\|. 
$$
Since 
$$
\int_I |\partial_x u|^2 - \lambda_1 |u|^2 = \int_I |\partial_x v|^2 - \lambda_1
|v|^2,
$$
from \eqref{eq:spectralgap1} we deduce the estimates
\begin{equation}\label{eq:controlv}
	\|v\|^2 \leq \frac{1}{\lambda_2 - \lambda_1}\|f\|\: \|u\|\quad \text{
	and } \quad
	\|\partial_x v\|^2 \leq \big(1 - \frac{\lambda_1}{\lambda_2}\big)^{-1}
\|f\|\: \|u\|.
\end{equation}
If instead we retain the imaginary part, we obtain
\begin{equation}\label{eq:controlweightedu}
\left| \int_I (V - s) |u|^2 \right| \leq \|f\|\: \|u\|.
\end{equation}

We shall distinguish two cases for $s$.

\smallskip

\noindent{\bf Case 1:} $\|V - s\|_\infty \geq Osc(V) + \omega_2(V).$\\
Then in \eqref{eq:controlweightedu} the 
weight $(V - s)$ is of constant sign a.e. on $I$ and moreover 
by definition of $Osc(V)$ 
$$
|V - s| \geq \omega_2(V) \qquad a.e. \text{ on } I.
$$
We therefore deduce from \eqref{eq:controlweightedu} that
\begin{equation}\label{eq:case1}
\|u \| \leq \frac{1}{\omega_2(V)} \|f\|.
\end{equation}

\medskip
\noindent{\bf Case 2:} $\|V - s\|_\infty \leq Osc(V) + \omega_2(V).$\\
In this case, the difficulty is that $V - s$ possibly changes sign. We aim to use the 
decomposition \eqref{eq:decomp-u} and to control $u$ by (controlling) 
$C_1=\langle u, e_1 \rangle$. To do so, we consider some $\phi \in Lip(I, \R)$ 
such that $\phi e_1(a) = \phi e_1(b)$,
multiply equation \eqref{eq:resolv} by $\phi \bar u$ and then retain the
imaginary part. Using that $e_1$ is real, this yields the identity 
\begin{equation}\label{eq:weightedident}\begin{split}
&{\rm Im}\Big(
C_1 \int_I \partial_x e_1 \bar v \phi' + \overline{C_1} \int_I e_1 \partial_x v \phi'
+ \int_I \bar v \partial_x v \phi' \Big) + |C_1|^2 \int_I (V - s)\phi e_1^2\\
& + \int_I (V - s) \phi \left[ |v|^2 + 2 {\rm Re}(C_1 e_1\bar v)\right]
= {\rm Im}\Big( \int_I f\phi \bar u\Big),
\end{split}\end{equation}
from which it follows that
\begin{equation}\label{eq:controlC1}\begin{split}
|C_1|^2 \Big| \int_I (V-s)\phi e_1^2 \Big| \leq 
& \ \|f\| \|u\| \|\phi\|_\infty\\
	& + |C_1| \|\phi'\|_\infty \Big[ \sqrt{\lambda_1}\|v\| + \|\partial_x v\| \Big] \\
& + \|\phi'\|_\infty \|\partial_x v\| \|v\|\\
&+ \|V - s\|_\infty \|\phi\|_\infty \Big[ 2 |C_1| \|v\| + \|v\|^2\Big],
\end{split}\end{equation}
where we used the Cauchy-Schwarz inequality with $\|e_1\|=1$ and 
$\|\partial_x e_1\|=\sqrt{\lambda_1}$. We now assume moreover that $\phi$ is a maximizer 
for $\omega_2(V)$ in 
\eqref{eq:defomega2}, that is $\int_I \phi e_1^2 = 0$,
$\|\phi\|_\infty \leq 1$, $\|\phi'\|_\infty  \leq 2\pi/|I|=\sqrt{\lambda_2}$ and 
$\int V \phi e_1^2 = \omega_2(V).$
In particular,
$$
\int_I (V-s)\phi e_1^2 = \int_I V \phi e_1^2 = \omega_2(V).
$$
Rewriting \eqref{eq:controlC1} taking into account the latter as well as 
\eqref{eq:controlv} and the assumption $\|V-s\|_\infty \leq Osc(V) +\omega_2(V)$ 
we then obtain 
\begin{equation}\label{eq:controlC1bis}\begin{split}
	|C_1|^2 \leq & \frac{1}{\omega_2(V)}\Bigg[ \Big( 1 + 
	\frac{\lambda_2}{\lambda_2-\lambda_1}\big(1 + \frac{(
	Osc(V) + \omega_2(V))}{\lambda_2}\big)\Big)
	\|f\| \|u\|\\
	& + \frac{\lambda_2}{\sqrt{\lambda_2-\lambda_1}}\Big(1 +
	\sqrt{\frac{\lambda_1}{\lambda_2}} + \frac{2(Osc(V) + \omega_2(V))}{\lambda_2}\Big)
	|C_1| \|f\|^\frac12 \|u\|^\frac12 \Bigg]. 
\end{split}
\end{equation} 
Since $\|u\|^2 = |C_1|^2 + \|v\|^2$, adding \eqref{eq:controlC1bis} with 
the first inequality in \eqref{eq:controlv}, and then replacing $|C_1|$ by $\|u\|$ on the
r.h.s. of the resulting inequality we finally obtain, after division by $\|u\|$, 
\begin{equation}\label{eq:quadratic0}
\|u\| \leq \alpha \|f\| + \beta \|f\|^\frac12 \|u\|^\frac12,
\end{equation}
where
$$
\alpha = \frac{1}{\omega_2(V)}\Big( 1 + 
    \frac{\lambda_2}{\lambda_2- \lambda_1}\Big) + \frac{Osc(V) + \omega_2(V)}{\omega_2(V)
    (\lambda_2 - \lambda_1)} +\frac{1}{\lambda_2 - \lambda_1} 
$$
and
$$
\beta = \frac{1}{\omega_2(V) \sqrt{\lambda_2-\lambda_1}} \Big(\lambda_2(1 +
	\sqrt{\tfrac{\lambda_1}{\lambda_2}})  + 2(Osc(V) + \omega_2(V))\Big).
$$
Analyzing \eqref{eq:quadratic0} as an inequality for a quadratic polynomial in
$\|u\|^\frac{1}{2}$, and computing its largest root, it immediately follows 
that  
\begin{equation}\label{eq:case2}
\|u\| \leq (4\alpha + \beta^2) \|f\|.
\end{equation}
In the remaining part of this proof we use the short hand notations $\omega$ for
$\omega_2(V)$ and $O$ for ${\it Osc}(V).$ In both boundary conditions 
cases\footnote{As a matter of fact, for the periodic case we have 
$\lambda_1=0$ and $\lambda_2=(2\pi/|I|)^2$, while for the Dirichlet case we 
have $\lambda_1=(\pi/|I|)^2$ and $\lambda_2=(2\pi/|I|)^2$.}, we have $\lambda_2 = (2\pi/|I|)^2$,
$\lambda_2 / (\lambda_2 - \lambda_1) \leq \frac{4}{3},$ and $\lambda_1 /
\lambda_2 \leq \frac{1}{4}$. In particular
$$
\alpha \leq \frac{7}{3 \omega} + \frac{|I|^2}{3\pi^2}\left(2 + 
\frac{O}{\omega}\right)
$$
and
$$
\beta \leq \frac{2\sqrt{3}\pi}{\omega|I|} + \frac{2|I|}{\sqrt{3}\pi}\left( 1 +
\frac{O}{\omega}\right),
$$
and therefore
\begin{equation}\label{eq:boundalphabeta}
4\alpha + \beta^2 \leq \frac{28}{3\omega} + \frac{4|I|^2}{3\pi^2}
\big(3 + 3\frac{O}{\omega} + (\frac{O}{\omega})^2\big)  
+ \frac{12\pi^2}{|I|^2\omega^2} + \frac{8}{\omega}(1 +  \frac{O}{\omega}).
\end{equation}
Note that from the definition \eqref{eq:defomega2} of $\omega$, and using
the fact that $e_1$ is $L^2$ normalized, it follows that in all
circumstances $\omega \leq \frac{1}{2}O.$ 
In \eqref{eq:boundalphabeta}, we may therefore estimate
$$
\frac{1}{\omega} \leq \frac{1}{2}\frac{\pi}{|I|\omega} \frac{|I|}{\pi}
\frac{O}{\omega} \leq \frac{1}{8} \frac{\pi^2}{|I|^2\omega^2} +
\frac{1}{2} \frac{|I|^2}{\pi^2}(\frac{O}{\omega})^2,
$$
$$
\big(3 + 3\frac{O}{\omega} + (\frac{O}{\omega})^2\big) \leq \frac{13}{4}(
\frac{O}{\omega})^2,
$$
and 
$$
\frac{8}{\omega}(1+ \frac{O}{\omega}) \leq 12 \frac{1}{\omega} \frac{O}{\omega} 
\leq 4\frac{\pi^2}{|I|^2\omega^2} + 9
\frac{|I|^2}{\pi^2} (\frac{O}{\omega})^2,
$$
which yields
\begin{equation}\label{eq:boundalphabeta2}
4\alpha + \beta^2 \leq (16 + \frac{7}{6}) \frac{\pi^2}{|I|^2\omega^2} +
18\frac{|I|^2}{\pi^2} (\frac{O}{\omega})^2 \leq \frac{18}{\omega^2} 
\left( \frac{\pi^2}{|I|^2} +
\frac{|I|^2}{\pi^2} O^2\right).
\end{equation}

We are now in position to end the proof in all cases. Indeed, note that in 
Case 1 estimate \eqref{eq:case1} we may transform as above  
$$
\frac{1}{\omega}  \leq \frac{1}{2}\frac{1}{\omega}\frac{O}{\omega} 
\leq 
\frac{1}{4}\frac{\pi^2}{|I|^2\omega^2} + 
\frac{1}{4}\frac{|I|^2}{\pi^2}(\frac{O}{\omega})^2
\leq
\frac{1}{4\omega^2} \left( \frac{\pi^2}{|I|^2} +
\frac{|I|^2}{\pi^2} O^2\right),
$$
so that combining \eqref{eq:case1} with \eqref{eq:case2} and
\eqref{eq:boundalphabeta2} the conclusion \eqref{eq:prop1} follows.  
\end{proof}

\end{document}
